\section{Cocontinuous functors and morphisms of topoi}
\label{section-cocontinuous-morphism-topoi}

\noindent
It is clear from the above that a cocontinuous functor $u$
gives a morphism of topoi in the same direction as $u$.
Thus this is in the opposite direction from the morphism of topoi
associated (under certain conditions) to a continuous $u$ as in
Definition \ref{definition-morphism-sites},
Proposition \ref{proposition-get-morphism}, and
Lemma \ref{lemma-morphism-sites-topoi}.

\begin{lemma}
\label{lemma-cocontinuous-morphism-topoi}
Let $\mathcal{C}$ and $\mathcal{D}$ be sites.
Let $u : \mathcal{C} \to \mathcal{D}$ be cocontinuous.
The functors $g_* = {}_su$ and $g^{-1} = (u^p\ )^\#$
define a morphism of topoi
$g$ from  $\Sh(\mathcal{C})$ to $\Sh(\mathcal{D})$.
\end{lemma}

\begin{proof}
This is exactly the content of Lemma \ref{lemma-exact-cocontinuous}.
\end{proof}

\begin{lemma}
\label{lemma-composition-cocontinuous}
\begin{slogan}
Composition of site functors respects cocontinuity of site functors.
\end{slogan}
Let $u : \mathcal{C} \to \mathcal{D}$, and $v : \mathcal{D} \to \mathcal{E}$
be cocontinuous functors. Then $v \circ u$ is cocontinuous and we
have $h = g \circ f$
where $f : \Sh(\mathcal{C}) \to \Sh(\mathcal{D})$,
resp.\ $g : \Sh(\mathcal{D}) \to \Sh(\mathcal{E})$,
resp.\ $h : \Sh(\mathcal{C}) \to \Sh(\mathcal{E})$ is the
morphism of topoi associated to $u$, resp.\ $v$, resp.\ $v \circ u$.
\end{lemma}

\begin{proof}
Let $U \in \Ob(\mathcal{C})$.
Let $\{E_i \to v(u(U))\}$ be a covering of $U$ in $\mathcal{E}$.
By assumption there exists a covering $\{D_j \to u(U)\}$ in $\mathcal{D}$
such that $\{v(D_j) \to v(u(U))\}$ refines $\{E_i \to v(u(U))\}$. Also
by assumption there exists a covering $\{C_l \to U\}$ in $\mathcal{C}$
such that $\{u(C_l) \to u(U)\}$ refines $\{D_j \to u(U)\}$. Then it is
true that $\{v(u(C_l)) \to v(u(U))\}$ refines the covering
$\{E_i \to v(u(U))\}$. This proves that $v \circ u$ is cocontinuous.
To prove the last assertion it suffices to show that
${}_sv \circ {}_su = {}_s(v \circ u)$. It suffices to prove that
${}_pv \circ {}_pu = {}_p(v \circ u)$, see Lemma \ref{lemma-pu-sheaf}.
Since ${}_pu$, resp.\ ${}_pv$, resp.\ ${}_p(v \circ u)$ is right adjoint to
$u^p$, resp.\ $v^p$, resp.\ $(v \circ u)^p$ it suffices to prove that
$u^p \circ v^p = (v \circ u)^p$. And this is direct from the definitions.
\end{proof}

\begin{example}
\label{example-open-immersion-cocontinuous}
Let $X$ be a topological space.
Let $j : U  \to X$ be the inclusion of an open subspace.
Recall that we have sites $X_{Zar}$ and $U_{Zar}$,
see Example \ref{example-site-topological}.
Recall that we have the functor $u : X_{Zar} \to U_{Zar}$
associated to $j$ which is continuous and gives rise
to a morphism of sites $U_{Zar} \to X_{Zar}$,
see Example \ref{example-continuous-map}.
This also gives a morphism of topoi $(j_*, j^{-1})$.
Next, consider the functor
$v : U_{Zar} \to X_{Zar}$, $V \mapsto v(V) = V$
(just the same open but now thought of as an object of $X_{Zar}$).
This functor is cocontinuous. Namely, if $v(V) = \bigcup_{j \in J} W_j$
is an open covering in $X$, then each $W_j$ must be a subset of
$U$ and hence is of the form $v(V_j)$,
and trivially $V = \bigcup_{j \in J} V_j$ is an open covering in $U$.
We conclude by Lemma \ref{lemma-cocontinuous-morphism-topoi} above
that there is a morphism of topoi associated to $v$
$$
\Sh(U) \longrightarrow \Sh(X)
$$
given by ${}_sv$ and $(v^p\ )^\#$. We claim that
actually $(v^p\ )^\# = j^{-1}$ and that ${}_sv = j_*$,
in other words, that this is the same morphism of topoi
as the one given above. Perhaps the easiest way to see
this is to realize that for any sheaf
$\mathcal{G}$ on $X$ we have
$v^p\mathcal{G}(V) = \mathcal{G}(V)$ which according to
Sheaves, Lemma \ref{sheaves-lemma-j-pullback} is a description
of $j^{-1}\mathcal{G}$ (and hence sheafification is superfluous
in this case). The equality of ${}_sv$ and $j_*$ follows by
uniqueness of adjoint functors (but may also be computed directly).
\end{example}

\begin{example}
\label{example-open-map-cocontinuous}
This example is a slight generalization of
Example \ref{example-open-immersion-cocontinuous}.
Let $f : X \to Y$ be a continuous map of topological spaces.
Assume that $f$ is open.
Recall that we have sites $X_{Zar}$ and $Y_{Zar}$,
see Example \ref{example-site-topological}.
Recall that we have the functor $u : Y_{Zar} \to X_{Zar}$
associated to $f$ which is continuous and gives rise
to a morphism of sites $X_{Zar} \to Y_{Zar}$,
see Example \ref{example-continuous-map}.
This also gives a morphism of topoi $(f_*, f^{-1})$.
Next, consider the functor
$v : X_{Zar} \to Y_{Zar}$, $U \mapsto v(U) = f(U)$.
This functor is cocontinuous. Namely, if $f(U) = \bigcup_{j \in J} V_j$
is an open covering in $Y$, then setting $U_j = f^{-1}(V_j) \cap U$
we get an open covering $U = \bigcup U_j$ such that $f(U) = \bigcup f(U_j)$
is a refinement of $f(U) = \bigcup V_j$.
We conclude by Lemma \ref{lemma-cocontinuous-morphism-topoi} above
that there is a morphism of topoi associated to $v$
$$
\Sh(X) \longrightarrow \Sh(Y)
$$
given by ${}_sv$ and $(v^p\ )^\#$. We claim that
actually $(v^p\ )^\# = f^{-1}$ and that ${}_sv = f_*$,
in other words, that this is the same morphism of topoi
as the one given above. For any sheaf
$\mathcal{G}$ on $Y$ we have
$v^p\mathcal{G}(U) = \mathcal{G}(f(U))$.
On the other hand, we may compute
$u_p\mathcal{G}(U) = \colim_{f(U) \subset V} \mathcal{G}(V)
= \mathcal{G}(f(U))$ because clearly $(f(U), U \subset f^{-1}(f(U)))$
is an initial object of the category $\mathcal{I}_U^u$ of
Section \ref{section-functoriality-PSh}.
Hence $u_p = v^p$ and we conclude $f^{-1} = u_s = (v^p\ )^\#$.
The equality of ${}_sv$ and $f_*$ follows by
uniqueness of adjoint functors (but may also be computed directly).
\end{example}

\noindent
In the first Example \ref{example-open-immersion-cocontinuous}
the functor $v$ is also continuous. But in the second
Example \ref{example-open-map-cocontinuous} it is generally not continuous
because condition (2) of Definition \ref{definition-continuous} may fail.
Hence the following lemma applies to the first example, but not to the second.

\begin{lemma}
\label{lemma-when-shriek}
\begin{slogan}
Sheafification is redundant in topoi morphisms associated to simultaneously
continuous and cocontinuous site functors.
\end{slogan}
Let $\mathcal{C}$ and $\mathcal{D}$ be sites.
Let $u : \mathcal{C} \to \mathcal{D}$ be a functor.
Assume that
\begin{enumerate}
\item[(a)] $u$ is cocontinuous, and
\item[(b)] $u$ is continuous.
\end{enumerate}
Let $g : \Sh(\mathcal{C}) \to \Sh(\mathcal{D})$
be the associated morphism of topoi. Then
\begin{enumerate}
\item sheafification in the formula $g^{-1} = (u^p\ )^\#$ is
unnecessary, in other words $g^{-1}(\mathcal{G})(U) = \mathcal{G}(u(U))$,
\item $g^{-1}$ has a left adjoint $g_{!} = (u_p\ )^\#$, and
\item $g^{-1}$ commutes with arbitrary limits and colimits.
\end{enumerate}
\end{lemma}

\begin{proof}
By Lemma \ref{lemma-pushforward-sheaf} for any sheaf $\mathcal{G}$
on $\mathcal{D}$ the presheaf $u^p\mathcal{G}$ is a sheaf on $\mathcal{C}$
(and is denoted $u^s\mathcal{G}$). The functor
$\mathcal{F} \mapsto (u_p\mathcal{F})^\#$ is a left
adjoint to $u^s$ by Lemma \ref{lemma-adjoint-sheaves}.
The statement on limits and colimits follows from the
discussion in Categories, Section \ref{categories-section-adjoint}.
\end{proof}

\noindent
In the situation of Lemma \ref{lemma-when-shriek} above we see that we have
a sequence of adjoint functors
$$
g_{!}, \ g^{-1}, \ g_*.
$$
The functor $g_!$ is {\it not} exact in general, because it does
not transform a final object of $\Sh(\mathcal{C})$
into a final object of $\Sh(\mathcal{D})$ in general.
See Sheaves, Remark \ref{sheaves-remark-j-shriek-not-exact}.
On the other hand, in the topological setting of
Example \ref{example-open-immersion-cocontinuous} the functor
$j_!$ is exact on abelian sheaves, see
Modules, Lemma \ref{modules-lemma-j-shriek-exact}. The following lemma gives
the generalization to the case of sites.

\begin{lemma}
\label{lemma-preserve-equalizers}
Let $\mathcal{C}$ and $\mathcal{D}$ be sites.
Let $u : \mathcal{C} \to \mathcal{D}$ be a functor.
Assume that
\begin{enumerate}
\item[(a)] $u$ is cocontinuous,
\item[(b)] $u$ is continuous, and
\item[(c)] fibre products and equalizers exist in $\mathcal{C}$ and
$u$ commutes with them.
\end{enumerate}
In this case the functor $g_!$ above commutes with fibre products and
equalizers (and more generally with finite connected limits).
\end{lemma}

\begin{proof}
Assume (a), (b), and (c).
We have $g_! = (u_p\ )^\#$. Recall (Lemma \ref{lemma-limit-sheaf}) that
limits of sheaves are equal to the corresponding limits as presheaves.
And sheafification commutes with finite limits
(Lemma \ref{lemma-sheafification-exact}). Thus it
suffices to show that $u_p$ commutes with fibre products and equalizers.
To do this it suffices that colimits over the categories
$(\mathcal{I}_V^u)^{opp}$ of
Section \ref{section-functoriality-PSh}
commute with fibre products and equalizers. This follows
from
Lemma \ref{lemma-almost-directed}
and
Categories, Lemma \ref{categories-lemma-almost-directed-commutes-equalizers}.
\end{proof}

\noindent
The following lemma deals with a case that is even more like
the morphism associated to an open immersion of topological spaces.

\begin{lemma}
\label{lemma-back-and-forth}
Let $\mathcal{C}$ and $\mathcal{D}$ be sites.
Let $u : \mathcal{C} \to \mathcal{D}$ be a functor.
Assume that
\begin{enumerate}
\item[(a)] $u$ is cocontinuous,
\item[(b)] $u$ is continuous, and
\item[(c)] $u$ is fully faithful.
\end{enumerate}
For $g_!, g^{-1}, g_*$ as above
the canonical maps $\mathcal{F} \to g^{-1}g_!\mathcal{F}$
and $g^{-1}g_*\mathcal{F} \to \mathcal{F}$ are isomorphisms
for all sheaves $\mathcal{F}$ on $\mathcal{C}$.
\end{lemma}

\begin{proof}
Let $X$ be an object of $\mathcal{C}$.
In Lemmas \ref{lemma-pu-sheaf} and \ref{lemma-when-shriek} we have seen that
sheafification is not necessary for the functors
$g^{-1} = (u^p\ )^\#$ and $g_{*} = ({}_pu\ )^\#$.
We may compute
$(g^{-1}g_{*}\mathcal{F})(X) = g_{*}\mathcal{F}(u(X))
= \lim \mathcal{F}(Y)$. Here the limit
is over the category of pairs $(Y, u(Y) \to u(X))$
where the morphisms $u(Y) \to u(X)$ are not required to be
of the form $u(\alpha)$ with $\alpha$ a morphism of $\mathcal{C}$.
By assumption (c) we see that they automatically come from
morphisms of $\mathcal{C}$ and we deduce that the limit is the
value on $(X, u(\text{id}_X))$, i.e., $\mathcal{F}(X)$.
This proves that $g^{-1}g_{*}\mathcal{F} = \mathcal{F}$.

\medskip\noindent
On the other hand, $(g^{-1}g_{!}\mathcal{F})(X) =
g_{!}\mathcal{F}(u(X)) = (u_p\mathcal{F})^\#(u(X))$, and
$u_p\mathcal{F}(u(X)) = \colim \mathcal{F}(Y)$.
Here the colimit is over the category of pairs $(Y, u(X) \to u(Y))$
where the morphisms $u(X) \to u(Y)$ are not required to be
of the form $u(\alpha)$ with $\alpha$ a morphism of $\mathcal{C}$.
By assumption (c) we see that they automatically come
from morphisms of $\mathcal{C}$ and we deduce that the colimit is the
value on $(X, u(\text{id}_X))$, i.e., $\mathcal{F}(X)$. Thus for every
$X \in \Ob(\mathcal{C})$ we have
$u_p\mathcal{F}(u(X)) = \mathcal{F}(X)$.
Since $u$ is cocontinuous and continuous any covering of $u(X)$ in
$\mathcal{D}$ can be refined by a covering (!) $\{u(X_i) \to u(X)\}$
of $\mathcal{D}$ where $\{X_i \to X\}$ is a covering in $\mathcal{C}$.
This implies that $(u_p\mathcal{F})^+(u(X)) = \mathcal{F}(X)$ also,
since in the colimit defining the value of $(u_p\mathcal{F})^+$
on $u(X)$ we may restrict to the cofinal system of coverings
$\{u(X_i) \to u(X)\}$ as above. Hence we see that
$(u_p\mathcal{F})^+(u(X)) = \mathcal{F}(X)$ for all objects $X$
of $\mathcal{C}$ as well. Repeating this argument one more time
gives the equality $(u_p\mathcal{F})^\#(u(X)) = \mathcal{F}(X)$
for all objects $X$ of $\mathcal{C}$. This produces the desired
equality $g^{-1}g_!\mathcal{F} = \mathcal{F}$.
\end{proof}

\noindent
Finally, here is a case that does not have any corresponding topological
example. We will use this lemma to see what happens when we enlarge a
``partial universe'' of schemes keeping the same topology. In the situation
of the lemma, the morphism of topoi $g : \Sh(\mathcal{C}) \to \Sh(\mathcal{D})$
identifies $\Sh(\mathcal{C})$ as a subtopos of $\Sh(\mathcal{D})$
(Section \ref{section-subtopoi}) and moreover, the given embedding has a
retraction.

\begin{lemma}
\label{lemma-bigger-site}
Let $\mathcal{C}$ and $\mathcal{D}$ be sites.
Let $u : \mathcal{C} \to \mathcal{D}$ be a functor.
Assume that
\begin{enumerate}
\item[(a)] $u$ is cocontinuous,
\item[(b)] $u$ is continuous,
\item[(c)] $u$ is fully faithful,
\item[(d)] fibre products exist in $\mathcal{C}$ and $u$ commutes with them,
and
\item[(e)] there exist final objects
$e_\mathcal{C} \in \Ob(\mathcal{C})$,
$e_\mathcal{D} \in \Ob(\mathcal{D})$ such that
$u(e_\mathcal{C}) = e_\mathcal{D}$.
\end{enumerate}
Let $g_!, g^{-1}, g_*$ be as above. Then, $u$ defines a morphism of sites
$f : \mathcal{D} \to \mathcal{C}$ with $f_* = g^{-1}$, $f^{-1} = g_!$.
The composition
$$
\xymatrix{
\Sh(\mathcal{C}) \ar[r]^g &
\Sh(\mathcal{D}) \ar[r]^f &
\Sh(\mathcal{C})
}
$$
is isomorphic to the identity morphism of the topos
$\Sh(\mathcal{C})$. Moreover, the functor $f^{-1}$ is fully faithful.
\end{lemma}

\begin{proof}
By assumption the functor $u$ satisfies the hypotheses of
Proposition \ref{proposition-get-morphism}. Hence $u$ defines
a morphism of sites and hence a morphism of topoi $f$ as in
Lemma \ref{lemma-morphism-sites-topoi}. The formulas
$f_* = g^{-1}$ and $f^{-1} = g_!$ are clear from the lemma cited and
Lemma \ref{lemma-when-shriek}.
We have
$f_* \circ g_* = g^{-1} \circ g_* \cong \text{id}$, and
$g^{-1} \circ f^{-1} = g^{-1} \circ g_! \cong \text{id}$
by Lemma \ref{lemma-back-and-forth}.

\medskip\noindent
We still have to show that $f^{-1}$ is fully faithful.
Let $\mathcal{F}, \mathcal{G} \in \Ob(\Sh(\mathcal{C}))$.
We have to show that the map
$$
\Mor_{\Sh(\mathcal{C})}(\mathcal{F}, \mathcal{G})
\longrightarrow
\Mor_{\Sh(\mathcal{D})}(f^{-1}\mathcal{F}, f^{-1}\mathcal{G})
$$
is bijective. But the right hand side is equal to
\begin{align*}
\Mor_{\Sh(\mathcal{D})}(f^{-1}\mathcal{F}, f^{-1}\mathcal{G})
& =
\Mor_{\Sh(\mathcal{C})}(\mathcal{F}, f_*f^{-1}\mathcal{G}) \\
& =
\Mor_{\Sh(\mathcal{C})}(\mathcal{F}, g^{-1}f^{-1}\mathcal{G}) \\
& =
\Mor_{\Sh(\mathcal{C})}(\mathcal{F}, \mathcal{G})
\end{align*}
(the first equality by adjunction) which proves what we want.
\end{proof}

\begin{example}
\label{example-closed-map-cocontinuous-false}
Let $X$ be a topological space.
Let $i : Z  \to X$ be the inclusion of a subset (with induced
topology). Consider the functor $u : X_{Zar} \to Z_{Zar}$,
$U \mapsto u(U) = Z \cap U$.
At first glance it may appear that this functor is cocontinuous as
well. After all, since $Z$ has the induced topology,
shouldn't any covering of $U\cap Z$ it come from a covering
of $U$ in $X$? Not so! Namely, what if $U \cap Z = \emptyset$?
In that case, the empty covering is a covering of
$U \cap Z$, and the empty covering can only be refined by
the empty covering. Thus we conclude that
$u$ cocontinuous $\Rightarrow$ every nonempty open $U$
of $X$ has nonempty intersection with $Z$.
But this is not sufficient. For example, if
$X = \mathbf{R}$ the real number line with the usual
topology, and $Z = \mathbf{R} \setminus \{0\}$, then
there is an open covering of $Z$, namely
$Z = \{x < 0\} \cup \bigcup_n \{1/n < x\}$ which
cannot be refined by the restriction of any
open covering of $X$.
\end{example}







